\mychapter{Continuidade e Validação Final}{cap:plano_continuidade}\lhead{CONTINUIDADE E VALIDAÇÃO FINAL}

A metodologia AIPE está formulada e possui duas aplicações experimentais. A Bahia e o piloto M5 têm matrizes completas. O próximo passo é construir a avaliação final do seletor, e não apresentar a rodada M5 interrompida como experimento concluído. O piloto já orientou diagnóstico e recuo; integra o conjunto de desenvolvimento. A expansão nacional permanece cancelada.

\section{Separar desenvolvimento e confirmação}
\label{sec:estudos_caso_adicionais}
A próxima amostra deve oferecer mais exemplos por classe e registrar suporte de cada perfil. A concentração natural do M5 em séries esparsas impede presumir cobertura equilibrada. Elegibilidade e estratificação usam somente o treino; séries reservadas não podem orientar o ajuste do seletor.

Hiperparâmetros, critérios de recuo e eventual simplificação de rótulos devem ser definidos em validação interna por loja. O teste final fica isolado. Uma melhora encontrada no piloto não é uma segunda replicação independente \citep{CawleyTalbot2010}. Ajustes de políticas compatíveis podem ser reaproveitados após auditoria de dados, configuração e código; mudanças científicas exigem novo ajuste da candidata afetada.

\section{Harmonizar o problema de decisão}
A calibração da CappedBaseStock precisa ser harmonizada entre bases antes de atribuir diferenças à demanda. Também é necessário comparar restrições de serviço com escopo claro: série, perfil ou rede. A meta de 0,70 e os custos atuais são parâmetros experimentais, não exigências empresariais já justificadas.

A comparação reduzida por consequências e recuo fornece um candidato para continuidade, com benefício concentrado em três intervenções. A instabilidade dos rótulos também motiva examinar conjuntos de candidatas quase empatadas e alvos sensíveis ao custo de erro. Essa análise deve ocorrer no desenvolvimento, sem selecionar a formulação pelo teste reservado. GA--DQN continua acompanhado por GA e DQN, com o esforço adicional informado.

\section{Sensibilidades e eficiência}
\label{sec:robustez_pooling}
Prazo, custos relativos, estoque inicial, tratamento terminal e horizonte podem alterar o ranking e a viabilidade física. Seu efeito deverá ser separado do efeito do seletor. Calibração agrupada por perfil é outra hipótese: não equivale a seleção de família por perfil e exige comparação própria.

A escolha de CPU ou GPU deve considerar tempo total sob a carga de uso prevista. O piloto utilizou oito processos CPU; essa configuração operacional não é uma vantagem científica de uma política. Não há execução nacional ou expansão automática autorizada como requisito de conclusão desta análise.

\section{Critérios de fechamento}
\begin{table}[htbp]\centering\small
\caption{Critérios de conclusão, sem datas ainda não acordadas.}\label{tab:cronograma}
\begin{tabular}{@{}p{4cm}p{8.8cm}@{}}\toprule
Etapa & Evidência necessária\\\midrule
Formulação do AIPE & Entradas, etapas, saídas e limites explicitados independentemente das bases\\
Protocolo final & Objetivo, suporte e calibrações fixados antes do teste\\
PSE & Comparação independente com global, perfil e controles simples\\
Robustez & Falhas físicas, portfólio, seleção e sensibilidades separadas\\
Redação & Afirmações alinhadas à evidência e revisão dos autores\\\bottomrule
\end{tabular}\end{table}

O artigo concentra a arquitetura e sua avaliação operacional; a dissertação mantém fundamentos, formulação e limites mais amplos. O material SBPO permanece arquivado. Autoria, declarações, disponibilidade de dados e elementos institucionais precisam ser fechados pelos responsáveis antes de submissão ou depósito.
